\clearpage
\section{Parallel argument filling and depth padding}
\label{sec:copying-parallel-padding}
This section proves the remaining individual constructor laws for actual
addressed filling: ordered parallel argument assembly, both arguments of
application and binary operations, and any finite amount of representation
padding. These are laws about the existing finite graphs and their actual
replacement operation. They do not change the term language, the proposition's
programs, the Figure~16 translation, or its variable/atom ordering.

The dependency chain is generic parallel replacement, recursive parallel
congruence, arbitrary-depth component addresses, the full pair and operation
clauses, and padding of addresses and fillings. The previous section supplies
abstraction and unary filling. The unrestricted context induction must still
combine these ingredients with the maximum-depth choices of the constructor
interpreter. In particular, a padded host requires padding the replacement at
the hole's local allowance, not merely at the host's allowance. Execution
invariance under these representation changes remains a separate obligation
of the copying machine. No observational equivalence or robustness claim is
inferred solely from a graph isomorphism.

\subsection{Replacement in an ordered parallel family}
\termdef{copying-parallel-filling}

Let $F_b$ be a finite family of graphs, let $a$ be a selected component,
and let $e$ be an edge of $F_a$. A graph $H$ has exactly the input and output
interface of $e$. The updated family replaces $e$ by $H$ in component $a$
and retains every other component. Its two evaluation lemmas split on whether
the component index equals $a$, eliminating the dependent interface transport
in the selected case. The parallel boundary is arbitrary but fixed; it may
permute component boundary blocks.


\lean{Graph/ParallelPlug}{15}{19}{ParallelPlug.family}

\lean{Graph/ParallelPlug}{20}{22}{ParallelPlug.family_same}

\lean{Graph/ParallelPlug}{23}{26}{ParallelPlug.family_other}

Before regrouping, internal vertices are either a tagged vertex of an old
component or a vertex of the inserted graph. After regrouping, they are tagged
vertices of the updated family. The forward map keeps every unchanged component
tag, marks old vertices in component $a$ as retained, and places inserted
vertices in component $a$ with the inserted tag. The two evaluation lemmas
identify its action on old vertices in the selected and unselected components.
The inverse inspects the component tag and, in component $a$, the retained or
inserted tag. Splitting on those same alternatives proves both inverse laws.
Thus the regrouping is a bijection without any assumptions about vertex names.


\lean{Graph/ParallelPlug}{27}{37}{ParallelPlug.vertexTo}

\lean{Graph/ParallelPlug}{38}{42}{ParallelPlug.vertexTo_same}

\lean{Graph/ParallelPlug}{43}{46}{ParallelPlug.vertexTo_other}

\lean{Graph/ParallelPlug}{47}{56}{ParallelPlug.vertexFrom}

\lean{Graph/ParallelPlug}{57}{80}{ParallelPlug.vertices}

The edge bijection uses the analogous regrouping, but the old carrier omits
the selected edge. A retained edge in component $a$ carries a proof that it is
not $e$; inequality of tagged edges supplies this proof. Edges of other
components keep their names and tags. New edges are placed in the selected
component. The inverse either returns the unchanged tagged edge or strips the
inserted tag. Case analysis proves the inverse identities, with equality
transport of carriers cancelling in each case.


\lean{Graph/ParallelPlug}{81}{91}{ParallelPlug.edgeTo}

\lean{Graph/ParallelPlug}{92}{104}{ParallelPlug.edgeFrom}

\lean{Graph/ParallelPlug}{105}{142}{ParallelPlug.edges}

Three attachment identities prove that regrouping commutes with the parallel
boundary: for a vertex of another component, a vertex of the selected component,
and a vertex inserted from $H$. Internal vertices follow the explicit maps.
Input and output vertices use the same boundary bijection on both sides.
For an inserted boundary vertex, attachment first reaches the selected edge's
source or target port, then the selected component's external embedding; this
is exactly its attachment after regrouping. The three edge evaluation lemmas
record the retained-selected, unselected, and newly inserted cases, allowing
the later dependent label and incidence expressions to be rewritten.


\lean{Graph/ParallelPlug}{143}{154}{ParallelPlug.map_other}

\lean{Graph/ParallelPlug}{155}{167}{ParallelPlug.map_component}

\lean{Graph/ParallelPlug}{168}{183}{ParallelPlug.map_inserted}

\lean{Graph/ParallelPlug}{184}{190}{ParallelPlug.edgeTo_same}

\lean{Graph/ParallelPlug}{191}{196}{ParallelPlug.edgeTo_other}

\lean{Graph/ParallelPlug}{197}{203}{ParallelPlug.edgeTo_inserted}

Replacing the tagged edge after parallel assembly is isomorphic to assembling
the updated family. Use the vertex and edge bijections above. Labels agree by
the selected/unselected family equations. For sources and targets, split each
edge into the same three cases. Mapping an ordered incidence list twice equals
mapping the composite; apply the appropriate attachment identity to every
entry. This proves equality of entire ordered lists, not merely membership or
unordered incidence. The theorem introduces no acyclicity or self-loop premise.


\lean{Graph/ParallelPlug}{204}{276}{ParallelPlug.iso}

Parallel assembly also transports component-wise isomorphisms whose edge
labels are related by a supplied relation. Both carrier maps retain the
component index and apply that component's bijection. Their inverses work
component-wise. Boundary embedding commutes with the vertex map: internal
vertices use the component bijection, while external vertices use the fixed
parallel boundary. Component label and incidence laws then prove the assembled
related-label isomorphism by mapping this identity over each incidence list.
Instantiating the label relation with recursive box isomorphism gives the
congruence used below.


\lean{Graph/ParallelIso}{12}{18}{ParallelIso.vertices}

\lean{Graph/ParallelIso}{19}{25}{ParallelIso.edges}

\lean{Graph/ParallelIso}{26}{30}{ParallelIso.map_embed}

\lean{Graph/ParallelIso}{31}{49}{ParallelIso.iso}

\subsection{Deep sites in either parallel component}

Specialise the family to two graphs using the existing Boolean component
index and an arbitrary ordered boundary. The net abbreviation denotes that
exact parallel constructor. A left or right edge retains its local name and
receives the corresponding tag; its label is unchanged. Specialising the
generic replacement isomorphism proves the two shallow filling laws. The
updated family agrees component-wise with the pair containing the replaced
graph, and equality transport yields the stated interfaces. Specialising
parallel related-label congruence gives recursive congruence for both graphs.


\lean{Copying/ParallelContextPlug}{13}{15}{ParallelContextPlug.net}

\lean{Copying/ParallelContextPlug}{16}{16}{ParallelContextPlug.leftEdge}

\lean{Copying/ParallelContextPlug}{17}{18}{ParallelContextPlug.rightEdge}

\lean{Copying/ParallelContextPlug}{19}{21}{ParallelContextPlug.leftEdge_label}

\lean{Copying/ParallelContextPlug}{22}{25}{ParallelContextPlug.rightEdge_label}

\lean{Copying/ParallelContextPlug}{26}{36}{ParallelContextPlug.plug_left}

\lean{Copying/ParallelContextPlug}{37}{47}{ParallelContextPlug.plug_right}

\lean{Copying/ParallelContextPlug}{48}{53}{ParallelContextPlug.congr}

Changing a box body in either component commutes with parallel assembly.
The outer carrier bijections retain the component tags and local elements.
In the selected component, split an edge into the chosen box edge and every
other edge: only the chosen label changes. In the opposite component every
label is unchanged. Box metadata fixes the external interface, so neither
source nor target incidence changes. The boundary embeddings agree on internal,
input and output vertices, proving both ordered incidence laws.


\lean{Copying/ParallelContextPlug}{54}{105}{ParallelContextPlug.replaceBox_left}

\lean{Copying/ParallelContextPlug}{106}{157}{ParallelContextPlug.replaceBox_right}

Lift an arbitrary-depth occurrence into either component. A shallow address
uses the tagged edge and the label equality to transport the local interfaces.
An address inside a box uses the tagged enclosing box and retains its body
address. Thus parallel assembly changes neither the local allowance nor the
hole interface. Filling a lifted shallow site uses the shallow replacement
law; filling a deep site uses the box-body interchange law with the already
filled inner body. This covers arbitrary finite nesting without flattening
boxes or identifying distinct equally labelled occurrences.


\lean{Copying/ParallelContextPlug}{158}{166}{ParallelContextPlug.leftSite}

\lean{Copying/ParallelContextPlug}{167}{175}{ParallelContextPlug.rightSite}

\lean{Copying/ParallelContextPlug}{176}{185}{ParallelContextPlug.plugAt_left}

\lean{Copying/ParallelContextPlug}{186}{195}{ParallelContextPlug.plugAt_right}

To preserve a sole-hole certificate, assume the fixed opposite component
is hole-free. Any other edge is either in the selected component, where
inequality of tags reduces to inequality of local edges, or in the fixed
component, where hole-freedom applies directly. A shallow site's marked hole
label is preserved. A deep site's inner certificate is retained and the same
other-edge argument applies outside its enclosing box. This certifies precisely
the addresses used in the filling laws.


\lean{Copying/ParallelContextPlug}{196}{205}{ParallelContextPlug.left_other_free}

\lean{Copying/ParallelContextPlug}{206}{215}{ParallelContextPlug.right_other_free}

\lean{Copying/ParallelContextPlug}{216}{227}{ParallelContextPlug.leftSite_sole}

\lean{Copying/ParallelContextPlug}{228}{238}{ParallelContextPlug.rightSite_sole}

\subsection{The unchanged pair, application and binary clauses}

The Figure~16 pair uses the prescribed middle-block exchange of the two
variable/atom interfaces. Instantiate the parallel laws with that exact
boundary and apply the existing outer-interface equality cast. This defines
both pair addresses. The host-cast filling law followed by parallel filling
proves both pair filling laws. Parallel recursive congruence lifts through
the same cast. The component sole-hole certificates also survive this cast,
giving both certified pair addresses.


\lean{Copying/BinaryContextPlug}{9}{15}{BinaryContextPlug.pairLeftSite}

\lean{Copying/BinaryContextPlug}{16}{21}{BinaryContextPlug.pairRightSite}

\lean{Copying/BinaryContextPlug}{22}{28}{BinaryContextPlug.pairLeft_filling}

\lean{Copying/BinaryContextPlug}{29}{36}{BinaryContextPlug.pairRight_filling}

\lean{Copying/BinaryContextPlug}{37}{40}{BinaryContextPlug.pair_congr}

\lean{Copying/BinaryContextPlug}{41}{45}{BinaryContextPlug.pairLeftSite_sole}

\lean{Copying/BinaryContextPlug}{46}{51}{BinaryContextPlug.pairRightSite_sole}

An application is the serial composite of the application operation, the
ordered pair of arguments, and the original environment distributor. To lift
an argument address, first lift into the pair, then into the right side of
the operation prefix, and finally into the left side of the distributor
suffix. Compose the two serial filling isomorphisms with pair filling and
serial congruence. This proves the left and right application laws separately.
The operation and distributor contain no holes, so the serial sole-hole laws
and pair certificates prove uniqueness of the selected occurrence. Finally,
recursive congruence of both arguments follows by pair congruence and two
applications of serial congruence. No distributor is removed or replaced.


\lean{Copying/BinaryContextPlug}{52}{57}{BinaryContextPlug.applicationLeftSite}

\lean{Copying/BinaryContextPlug}{58}{65}{BinaryContextPlug.applicationLeft_filling}

\lean{Copying/BinaryContextPlug}{66}{73}{BinaryContextPlug.applicationLeftSite_sole}

\lean{Copying/BinaryContextPlug}{74}{79}{BinaryContextPlug.applicationRightSite}

\lean{Copying/BinaryContextPlug}{80}{87}{BinaryContextPlug.applicationRight_filling}

\lean{Copying/BinaryContextPlug}{88}{95}{BinaryContextPlug.applicationRightSite_sole}

\lean{Copying/BinaryContextPlug}{96}{100}{BinaryContextPlug.application_congr}

For each binary operation, use its original primitive operation edge in the
same serial construction. The left and right site definitions follow the
corresponding argument tag through the pair, prefix and suffix. Their filling
proofs compose serial and pair filling as above; their sole-hole proofs use
hole-freedom of this binary operation and the environment distributor. Recursive
congruence follows from pair and serial congruence. The operation parameter is
arbitrary, so these results do not restrict the surrounding operation alphabet.


\lean{Copying/BinaryContextPlug}{101}{106}{BinaryContextPlug.binaryLeftSite}

\lean{Copying/BinaryContextPlug}{107}{114}{BinaryContextPlug.binaryLeft_filling}

\lean{Copying/BinaryContextPlug}{115}{122}{BinaryContextPlug.binaryLeftSite_sole}

\lean{Copying/BinaryContextPlug}{123}{128}{BinaryContextPlug.binaryRightSite}

\lean{Copying/BinaryContextPlug}{129}{136}{BinaryContextPlug.binaryRight_filling}

\lean{Copying/BinaryContextPlug}{137}{144}{BinaryContextPlug.binaryRightSite_sole}

\lean{Copying/BinaryContextPlug}{145}{149}{BinaryContextPlug.binary_congr}

\subsection{Padding the address and the replacement together}
\termdef{copying-padded-filling}

Two equality transports are needed. Changing the local interface of a site
is equivalent, upon filling, to reindexing the replacement by the inverse
interface equalities. Transporting a site's host along graph equality leaves
the resulting filling equal to the old filling. Define the host transport by
equality elimination. Both filling equalities follow by eliminating the
relevant equalities, after which the reflexive reindexings disappear.


\lean{Copying/PaddingPlug}{9}{15}{PaddingPlug.plugAt_castInterfaces}

\lean{Copying/PaddingPlug}{16}{18}{PaddingPlug.Site.castGraph}

\lean{Copying/PaddingPlug}{19}{24}{PaddingPlug.plugAt_castGraph}

The uniform padding label map preserves primitive labels, and maps a box to
the same bound/free-variable/free-atom metadata with its body raised once.
At allowance zero it removes the primitive label's type-level wrapper. Splitting
on the depth and label shows that both ordered signatures are unchanged.
Relabelling by this map gives a graph with definitionally unchanged outer
carriers and incidence. It equals the existing recursive raise operation:
at depth zero the definitions coincide; at positive depth equality of the
label functions follows by primitive/box case analysis. This is a proof about
the existing padding, not a replacement definition of translation.


\lean{Copying/PaddingPlug}{25}{29}{PaddingPlug.label}

\lean{Copying/PaddingPlug}{30}{35}{PaddingPlug.label_input}

\lean{Copying/PaddingPlug}{36}{42}{PaddingPlug.label_output}

\lean{Copying/PaddingPlug}{43}{45}{PaddingPlug.graph}

\lean{Copying/PaddingPlug}{46}{55}{PaddingPlug.graph_eq_raise}

One-step padding commutes with shallow replacement up to ordinary isomorphism.
Both graphs retain the same old and inserted carrier tags, so identity carrier
bijections suffice. The labels coincide under the uniform padding map. For
retained edges the source and target lists coincide. For inserted edges, the
local signature casts cancel and each internal/input/output attachment agrees;
map composition proves equality of the ordered incidence lists.

Padding also commutes with changing a particular box body. Keep all outer
carrier names fixed. At the selected edge both labels are the same padded box
with the new padded body; all other labels agree. The outer incidence never
changes. This second law handles a site's enclosing box.


\lean{Copying/PaddingPlug}{56}{98}{PaddingPlug.plug}

\lean{Copying/PaddingPlug}{99}{117}{PaddingPlug.replaceBox_iso}

Raise an address by structural induction. A shallow site in the uniform
relabelling has the old local interface by the two signature equalities;
transport its host to the original raise operation. A deep site retains its
outer edge, pads the enclosing box body, and recursively raises its inner
address. Both the host and the local hole allowance increase by one.

The filling theorem states that padding the filled graph is recursively
isomorphic to filling the raised address with the raised replacement. In the
shallow case, the equality-transport laws reduce the statement to shallow
padding/replacement interchange. In the deep case, compose box-body interchange
with recursive congruence of the changed box body, using the induction
hypothesis for the inner address. The replacement is raised at its own local
depth, which may differ from the host depth.


\lean{Copying/PaddingPlug}{118}{128}{PaddingPlug.site}

\lean{Copying/PaddingPlug}{129}{144}{PaddingPlug.filling}

Host equality transport preserves a sole-hole certificate by equality
elimination. For every other outer edge, the existing per-label padding law
preserves hole-freedom; the uniform presentation has the same property by
primitive/box case analysis. Induction on the site now proves that its raised
address is still sole. A shallow marked hole stays the same primitive hole,
and every other edge stays hole-free. A deep address uses the inner induction
hypothesis and the same other-edge argument for its enclosing graph.


\lean{Copying/PaddingPlug}{145}{149}{PaddingPlug.site_castGraph_sole}

\lean{Copying/PaddingPlug}{150}{161}{PaddingPlug.graph_other_free}

\lean{Copying/PaddingPlug}{162}{181}{PaddingPlug.site_sole}

Padding preserves recursive isomorphism. Retain the original outer carrier
bijections and incidence equations. At depth zero equal primitive labels
remain equal. At positive depth, related primitives remain equal, mismatched
primitive/box cases are impossible, and related box bodies remain related by
the depth induction hypothesis. Bound and free metadata stay identical.

Iterate the raised-address construction $n$ times. A site with host allowance
$d$ and local allowance $r$ then has allowances $d+n$ and $r+n$. Iterating the
preceding congruence gives isomorphism preservation under the original padding
function. Induction on $n$ proves the filling law: raise the preceding filling
isomorphism, then compose it with the one-step filling theorem for the current
padded site and replacement. A separate induction iterates sole-hole
preservation. These statements retain the exact same number of padding steps
on the host and replacement.


\Needspace{8\baselineskip}
\lean{Copying/PaddingPlug}{182}{218}{PaddingPlug.congr}

\lean{Copying/PaddingPlug}{219}{224}{PaddingPlug.padSite}

\lean{Copying/PaddingPlug}{225}{232}{PaddingPlug.pad_congr}

\lean{Copying/PaddingPlug}{233}{241}{PaddingPlug.pad_filling}

\lean{Copying/PaddingPlug}{242}{249}{PaddingPlug.padSite_sole}

\subsection{Combined regressions and completion boundary}

The regressions combine two padding steps, an arbitrary-depth occurrence and
the full application clause, on each side separately. The fixed opposite
argument is any hole-free graph at the resulting allowance and may contain
boxes. The same explicitly constructed address is proved sole, and its filling
is recursively isomorphic to the application of the padded original filling.
The proof composes application filling with the inverse iterated-padding
filling law under application congruence. Thus both component tags, the ordered
boundary exchange, the serial distributors, and distinct host/local allowances
are checked together. These are structural regressions, not executions of a
copying machine; callback divergence and store behaviour still require item~2.

The individual parallel and padding filling laws are now discharged. Remaining
item~3 work is the unrestricted context induction, including cumulative padding
at maximum-depth joins, and the passage from structural correspondence to
observational congruence with proved machine representation laws. The other
Appendix~B.1 obligations are unchanged.


\lean{Copying/ParallelFillingExamples}{10}{24}{ParallelFillingExamples.paddedApplicationLeft_filling}

\lean{Copying/ParallelFillingExamples}{25}{37}{ParallelFillingExamples.paddedApplicationRight_filling}
